% ============================================================
% PORTADA
% ============================================================
{
  \setbeamertemplate{navigation symbols}{}
  \setbeamertemplate{footline}{}
  \begin{frame}[plain]
    \titlepage
  \end{frame}
}

% ============================================================
% TABLA DE CONTENIDOS
% ============================================================
{
  \setbeamertemplate{navigation symbols}{}
  \setbeamertemplate{footline}{}
  \begin{frame}{¿Qué veremos hoy?}
    \tableofcontents
  \end{frame}
}

% ============================================================
\section{Calentando motores}
% ============================================================

% ------------------------------------------------------------
\subsection{UIP vs.\ CIP}
% ------------------------------------------------------------

\begin{frame}{Motivación — ¿por qué nos importa?}

  En el ejercicio 2 derivamos la condición de Euler para bonos del país 1:
  \[
    \mathbb{E}_t[\hat{\lambda}_{1,t+1}] - \hat{\lambda}_{1,t}
    = Q_t + \hat{q}^a_{1,t} - \mathbb{E}_t[\hat{q}^a_{1,t+1}] - \hat{\beta}_{1,t}.
  \]

  \bigskip
  \textbf{Pregunta natural:} ¿qué tan lejos estamos de la paridad de tasas de interés
  estándar que conocemos del curso de macro? \\[0.5em]

  Para responderla necesitamos distinguir dos conceptos relacionados pero distintos:

  \bigskip
  \begin{columns}
    \column{0.48\textwidth}
    \centering
    \textcolor{peacockblue}{\textbf{CIP}} \\
    \small Covered Interest Parity \\
    \footnotesize (paridad \emph{cubierta})

    \column{0.48\textwidth}
    \centering
    \textcolor{spicy}{\textbf{UIP}} \\
    \small Uncovered Interest Parity \\
    \footnotesize (paridad \emph{descubierta})
  \end{columns}

\end{frame}

% ------------------------------------------------------------

\begin{frame}{CIP — Paridad Cubierta de Tasas de Interés}

  \begin{bluebox}{Covered Interest Parity (CIP)}
    La CIP establece que, en ausencia de fricciones y costos de transacción, el
    diferencial de tasas de interés entre dos países debe igualarse al premio forward
    del tipo de cambio:
    \[
      i_t - i^*_t = f_t - e_t,
    \]
    donde $i_t$ es la tasa doméstica, $i^*_t$ la extranjera, $e_t = \ln \mathcal{E}_t$
    el (log del) tipo de cambio spot y $f_t = \ln \mathcal{F}_t$ el forward.
  \end{bluebox}

  \bigskip
  \begin{block}{Intuición}
    Un inversor que toma una posición \textbf{cubierta} no enfrenta riesgo cambiario:
    vende divisa extranjera en el mercado forward al mismo tiempo que invierte.
    La CIP es una condición de \textbf{no-arbitraje puro} — se cumple casi siempre
    en datos de alta frecuencia.
  \end{block}

\end{frame}

% ------------------------------------------------------------

\begin{frame}{UIP — Paridad Descubierta de Tasas de Interés}

  \begin{bluebox}{Uncovered Interest Parity (UIP)}
    La UIP establece que, bajo neutralidad al riesgo, el diferencial de tasas debe
    igualarse a la \emph{variación esperada} del tipo de cambio:
    \[
      i_t - i^*_t = \mathbb{E}_t[e_{t+1}] - e_t = \mathbb{E}_t[\Delta e_{t+1}].
    \]
    El inversor \textbf{no} cubre su posición: acepta el riesgo cambiario.
  \end{bluebox}

  \bigskip
  \begin{block}{Diferencia clave con CIP}
    \begin{itemize}
      \item \textbf{CIP}: usa el tipo de cambio \emph{forward} $f_t$ (conocido en $t$).
            No hay incertidumbre. Se cumple empíricamente.
      \item \textbf{UIP}: usa la \emph{expectativa} $\mathbb{E}_t[e_{t+1}]$ (no observable).
            Requiere neutralidad al riesgo. \textbf{Falla sistemáticamente} en los datos
            (puzzle de la prima forward).
    \end{itemize}
  \end{block}

\end{frame}

% ------------------------------------------------------------

\begin{frame}{CIP vs.\ UIP — Intuición gráfica}

  \begin{columns}
    \column{0.5\textwidth}
    \textbf{CIP} \\[0.3em]
    \small
    Inversor doméstico: \\
    \quad 1. Invierte \$1 en bono doméstico: obtiene $(1+i_t)$ en $t+1$. \\
    \quad 2. Convierte \$1 a divisa, invierte en bono extranjero y \textbf{vende} divisa
             a plazo: obtiene $(1+i^*_t)\cdot\mathcal{F}_t/\mathcal{E}_t$.
    \\[0.5em]
    No-arbitraje: $(1+i_t) = (1+i^*_t)\cdot\mathcal{F}_t/\mathcal{E}_t$. \\
    En logs: \textcolor{peacockblue}{$i_t - i^*_t \approx f_t - e_t$}. \checkmark

    \column{0.5\textwidth}
    \textbf{UIP} \\[0.3em]
    \small
    Inversor doméstico: \\
    \quad 1. Invierte \$1 en bono doméstico: obtiene $(1+i_t)$.\\
    \quad 2. Convierte \$1 a divisa, invierte en bono extranjero y \textbf{no cubre}:
             obtiene $(1+i^*_t)\cdot\mathcal{E}_{t+1}/\mathcal{E}_t$.
    \\[0.5em]
    Indiferencia en esperanza: \\
    \textcolor{spicy}{$i_t - i^*_t \approx \mathbb{E}_t[\Delta e_{t+1}]$}. \\[0.3em]
    Requiere neutralidad al riesgo.
  \end{columns}

\end{frame}

% ------------------------------------------------------------

\begin{frame}{Conexión con el modelo de Corsetti \& Müller (2006)}

  La condición de Euler del modelo (pregunta 2) puede reescribirse como:
  \[
    Q_t = \underbrace{\mathbb{E}_t[\hat{q}^a_{1,t+1}] - \hat{q}^a_{1,t}}_{\text{análogo de } \mathbb{E}_t[\Delta e_{t+1}]}
    +\; \underbrace{\mathbb{E}_t[\hat{\lambda}_{1,t+1}] - \hat{\lambda}_{1,t}
    + \hat{\beta}_{1,t}}_{\text{prima de riesgo endógena}}.
  \]

  \bigskip
  \begin{block}{Lectura en clave UIP}
    \begin{itemize}
      \item El precio del bono $Q_t$ juega el papel del diferencial de tasas $i_t - i^*_t$.
      \item La variación esperada de $\hat{q}^a_1$ es el análogo de $\mathbb{E}_t[\Delta e]$:
            cuánto se deprecia el precio del bien en que está denominado el bono.
      \item Los términos de utilidad marginal y $\hat{\beta}_{1,t}$ son la \textbf{prima de
            riesgo}: desaparecen bajo neutralidad al riesgo o mercados completos.
      \item Por eso el ejercicio 2 es la base para entender la prima de riesgo cambiaria
            de Engel (1992) en el ejercicio 6.
    \end{itemize}
  \end{block}

\end{frame}

% ============================================================
\subsection{La historia de la replicación}
% ============================================================

\begin{frame}{La historia detrás del código}

  \textbf{¿De dónde viene el archivo \texttt{CM\_2006\_Klein.m}?}

  \bigskip
  \begin{enumerate}
    \item \textbf{Codeo manual a partir del paper.} Leímos las ecuaciones del modelo en
          \citet{corsetti2006twin} e intentamos transcribirlas directamente a MATLAB,
          siguiendo la convención $A\,\mathbb{E}_t[x_{t+1}] = B\,x_t$.

    \item \textbf{Lectura de papers técnicos.} Para entender el solver, revisamos
          \citet{klein2000} sobre la descomposición de Schur generalizada.

    \item \textbf{Uso de Claude.} Usamos un asistente de IA para depurar los índices de
          variables y las matrices $A$ y $B$ — \textit{error tras error}.

    \item \textbf{WayBack Machine.} Los códigos originales del Prof.\ Klein, P.\ habían
          sido removidos de su sitio web. Los recuperamos gracias al
          \href{https://web.archive.org}{\texttt{web.archive.org}} (WayBack Machine).

    \item \textbf{Escribir al Prof.\ Müller.} Cuando todos los diagnósticos fallaron,
          decidimos escribirle directamente al autor.
  \end{enumerate}

\end{frame}

% ------------------------------------------------------------

\begin{frame}{La respuesta del Profesor Müller}

  Después de describir nuestros problemas con el eigenvalor explosivo y las IRF con
  spikes, el Prof.\ Gernot Müller (Universidad de Tübingen) respondió en menos de 24h
  con sus códigos originales.

  \bigskip
  \begin{center}
    % -------------------------------------------------------
    % INSTRUCCIÓN: coloque aquí la captura de pantalla del
    % correo del Prof. Müller. Nombre sugerido: email_muller.png
    % -------------------------------------------------------
    \includegraphics[width=0.82\textwidth]{email_muller.png}
  \end{center}

\end{frame}

% ------------------------------------------------------------

\begin{frame}{¿Qué encontramos con los códigos originales?}

  Al comparar nuestro código con el \texttt{model\_IRF.m} del Prof.\ Müller, identificamos
  \textbf{cuatro errores estructurales}:

  \bigskip
  \begin{redbox}{Errores estructurales}
    \begin{enumerate}
      \item \textbf{FOC de capital, país 2:} usábamos $\hat{q}^a_{2,t}$ en vez de
            $\hat{q}^b_{2,t}$ (el país 2 produce el bien $b$, no el $a$).
      \item \textbf{Euler de bonos, país 2:} precio incorrecto en período corriente
            ($\hat{q}^b_{2,t}$ en vez de $\hat{q}^a_{2,t}$).
      \item \textbf{Clearing de bonos:} teníamos $b_1 = b_2$ en vez de $b_1 + b_2 = 0$
            (oferta neta nula).
      \item \textbf{Restricción presupuestaria del país 2:} faltaba completamente.
    \end{enumerate}
  \end{redbox}

  \bigskip
  \small Con los cuatro errores corregidos: eigenvalor máximo $= 0.9989$, IRF suaves,
  Figura~1 del paper replicada \textbf{idénticamente}. \happy

\end{frame}

% ============================================================
\subsection{Condiciones de Klein (2000)}
% ============================================================

\begin{frame}{El sistema de expectativas racionales lineales}

  \begin{bluebox}{Sistema de Klein (2000)}
    El modelo DSGE log-linealizado se escribe como:
    \[
      A\,\mathbb{E}_t[x_{t+1}] = B\,x_t,
    \]
    donde $x_t \in \mathbb{R}^n$ agrupa \textbf{todas} las variables del modelo.
    Las matrices $A, B \in \mathbb{R}^{n\times n}$ contienen los coeficientes
    de las ecuaciones log-linealizadas.
  \end{bluebox}

  \bigskip
  Particionamos $x_t$ en dos bloques:
  \[
    x_t = \begin{pmatrix} x^p_t \\ x^f_t \end{pmatrix},
  \]
  \begin{itemize}
    \item $x^p_t \in \mathbb{R}^{n_k}$: variables \textbf{predeterminadas} (estados),
          conocidas en $t$ (capital, bonos, shocks AR(1)).
    \item $x^f_t \in \mathbb{R}^{n-n_k}$: variables \textbf{no predeterminadas} (salto),
          determinadas en el equilibrio por las expectativas (consumo, precios, trabajo).
  \end{itemize}

  \bigskip
  En nuestro código: $n = 45$, $n_k = 9$ (ver índices $z_1,\ldots,b_1$).

\end{frame}

% ------------------------------------------------------------

\begin{frame}{Descomposición de Schur Generalizada}

  \begin{greenbox}{Teorema de Schur Generalizado}
    Para cualquier par de matrices $(A, B)$, existen matrices unitarias $Q, Z$ y
    matrices cuasi-triangulares superiores $S, T$ tales que:
    \[
      A = Q S Z^*, \qquad B = Q T Z^*.
    \]
    Los valores propios generalizados son los cocientes $\lambda_i = T_{ii}/S_{ii}$.
  \end{greenbox}

  \bigskip
  \begin{block}{¿Para qué sirve?}
    \begin{itemize}
      \item Cada $\lambda_i$ es un eigenvalor de la dinámica del sistema.
      \item $|\lambda_i| < 1$: modo \textbf{estable} (converge al SS).
      \item $|\lambda_i| > 1$: modo \textbf{inestable} (diverge).
      \item Ordenamos: primero los $n_k$ eigenvalores estables, luego los inestables.
    \end{itemize}
  \end{block}

  \bigskip
  \small \texttt{solab.m} implementa esto vía \texttt{qz(A,B)} de MATLAB y reordena
  automáticamente por módulo.

\end{frame}

% ------------------------------------------------------------

\begin{frame}{Condiciones de Blanchard-Kahn}

  \begin{greenbox}{Condiciones de existencia y unicidad (Blanchard-Kahn)}
    El sistema $A\,\mathbb{E}_t[x_{t+1}] = B\,x_t$ tiene una única solución estacionaria si y
    solo si el número de eigenvalores generalizados \emph{estables} $|\lambda_i| < 1$
    es exactamente igual al número de variables predeterminadas $n_k$:
    \[
      \#\{i : |\lambda_i| < 1\} = n_k.
    \]
  \end{greenbox}

  \bigskip
  \begin{columns}
    \column{0.48\textwidth}
    \textbf{Demasiados estables} ($> n_k$): \\
    \small indeterminación — múltiples equilibrios (burbujas).

    \column{0.48\textwidth}
    \textbf{Demasiados inestables} ($> n - n_k$): \\
    \small no existe solución estacionaria (Blanchard-Kahn falla). \skull
  \end{columns}

  \bigskip
  \begin{block}{En nuestro modelo}
    Con $n_k = 9$ estados, necesitamos exactamente 9 eigenvalores estables.
    Tras corregir los errores: máx$|\lambda_i| = 0.9989 < 1$ para los 9 modos estables.
    \checkmark
  \end{block}

\end{frame}

% ------------------------------------------------------------

\begin{frame}{La solución: \texttt{solab.m}}

  Una vez verificadas las condiciones de Blanchard-Kahn, la solución del sistema es
  \textbf{lineal} en los estados:
  \[
    x^f_t = F\, x^p_t, \qquad x^p_{t+1} = P\, x^p_t + \text{(shocks)},
  \]
  donde $F \in \mathbb{R}^{(n-n_k)\times n_k}$ es la \textbf{función de política} y
  $P \in \mathbb{R}^{n_k \times n_k}$ es la \textbf{ley de movimiento de los estados}.

  \bigskip
  \begin{block}{Uso en MATLAB}
    \begin{semiverbatim}
      [f\_sol, p\_sol] = solab(A, B, nk);
    \end{semiverbatim}
    \begin{itemize}
      \item \texttt{f\_sol} $\leftarrow$ $F$: dadas los estados hoy, predice todas las
            variables de salto.
      \item \texttt{p\_sol} $\leftarrow$ $P$: itera los estados hacia adelante.
    \end{itemize}
  \end{block}

  \bigskip
  \textbf{IRF:} fijamos el vector inicial $x^p_0 = e_{g_1}$ (shock unitario en $g_1$) y
  calculamos $x^p_t = P^{t-1} x^p_0$ y $x^f_t = F\,x^p_t$ para $t = 1,\ldots,40$.

\end{frame}

% ============================================================
% BIBLIOGRAFÍA
% ============================================================
\begin{frame}[allowframebreaks]{Referencias}
  \nocite{*}
  \printbibliography
\end{frame}
